% Módulo central mostrando la vista de configuración inicial (formulario).
% Se compila por separado: lualatex onx_config_formulario.tex -> onx_config_formulario.pdf
% Se dibuja sobre red_onx_apagada.png. La disposición replica screen_setup.cpp:
% título en Montserrat 18, etiquetas en Montserrat 14, campos de 280 px de ancho
% en y = 60, 145 y 230 px, y botón de 200x40 px en y = 310 px. El teclado es el
% de LVGL en modo minúsculas y ocupa la mitad inferior. Al enfocar la clave de
% tutor, el formulario sube 52 px para que el campo quede visible.
% Colores del firmware con la inversión del panel desactivada: fondo 0x2C3E50,
% etiquetas blancas, campos y teclado del tema claro de LVGL y botón 0x2196F3.
\documentclass[tikz,border=0pt]{standalone}
\usepackage{fontspec}
\usepackage{fontawesome5}
\newfontfamily\montserrat{Montserrat-Medium.otf}
\usetikzlibrary{arrows.meta}

\definecolor{pantalla}{RGB}{44,62,80}
\definecolor{oscuro}{RGB}{255,255,255}
\definecolor{campo}{RGB}{255,255,255}
\definecolor{placeholder}{RGB}{176,176,176}
\definecolor{icono}{RGB}{85,85,85}
\definecolor{claro}{RGB}{33,33,33}
\definecolor{naranja}{RGB}{33,150,243}
\definecolor{textoboton}{RGB}{255,255,255}
\definecolor{fondoteclado}{RGB}{230,232,235}
\definecolor{tecla}{RGB}{255,255,255}
\definecolor{teclaespecial}{RGB}{200,203,208}

\newcommand\fuentedieciocho{\montserrat\fontsize{8.2}{9.6}\selectfont}
\newcommand\fuentedieciseis{\montserrat\fontsize{7.3}{8.6}\selectfont}
\newcommand\fuentecatorce{\montserrat\fontsize{6.4}{7.6}\selectfont}

\begin{document}
\begin{tikzpicture}[x=0.01cm,y=0.01cm]
\node[anchor=south west, inner sep=0] at (0,0) {\includegraphics[width=6.39cm]{red_onx_apagada.png}};
% Área visible del panel en la imagen: (64,101)-(575,867), equivale a 320x480 px
\begin{scope}[shift={(64,867)}, x=0.015969cm, y=-0.015958cm]
  \fill[pantalla] (0,0) rectangle (320,480);
  \clip (0,0) rectangle (320,480);
  \begin{scope}[every node/.style={inner sep=0}]
    \node[anchor=north, font=\fuentedieciocho, text=oscuro] at (160,5) {L.E.V.I. - Configurar Sistema};
    \node[anchor=north, font=\fuentecatorce, text=oscuro] at (160,40) {Nombre del Alumno:};
    \fill[campo, rounded corners=0.06cm] (20,60) rectangle (300,112);
    \node[anchor=base west, font=\fuentedieciocho, text=placeholder] at (34,90.5) {Ej: Juan Perez};
    \node[anchor=north, font=\fuentecatorce, text=oscuro] at (160,125) {Clave Profesor:};
    \fill[campo, rounded corners=0.06cm] (20,145) rectangle (300,197);
    \node[anchor=base west, font=\fuentedieciocho, text=placeholder] at (34,175.5) {Clave Profesor};
    \node[font=\fontsize{7.3}{8}\selectfont, text=icono] at (279,171) {\faEyeSlash};
    \node[anchor=north, font=\fuentecatorce, text=oscuro] at (160,210) {Clave Tutor:};
    \fill[campo, rounded corners=0.06cm] (20,230) rectangle (300,282);
    \node[anchor=base west, font=\fuentedieciocho, text=placeholder] at (34,260.5) {Clave Tutor};
    \node[font=\fontsize{7.3}{8}\selectfont, text=icono] at (279,256) {\faEyeSlash};
    \fill[naranja, rounded corners=0.06cm] (60,310) rectangle (260,350);
    \node[font=\fuentecatorce, text=textoboton] at (160,330) {Guardar};
  \end{scope}
\end{scope}
\end{tikzpicture}
\end{document}
